%% First-order completeness, compactness, and cardinality.
%% Migrated from v1.1 Part III, lines 4306--5125.  Statements only;
%% proofs and the Henkin construction are in the parallel annex fragment.

\chapter{Completeness, Compactness, Cardinality}
\label{sec:p02-first-order-completeness}

For completeness and Henkin's construction, see
\cite{Godel1930,Henkin1949}.  We retain the source's set-based,
nonempty-domain semantics and exact choice assumptions; compare
\cite{Hodges1993,Marker2002}.

\section{Maximal consistency and witnesses}
\label{sec:p02-first-order-henkin-preparation}

\begin{hbdefinition}{def:p02-first-order-consistency}{First-order consistency}
For \(\Gamma\subseteq\operatorname{Fm}(\Lang)\), \(\Gamma\) is
\emph{consistent} iff
\(\Gamma\notderives_{\mathsf{ND}_{\mathsf{FOL}}}\bottom\).  It is
\emph{maximal consistent} iff it is consistent and
\(\Gamma\cup\{\varphi\}\) is inconsistent for every
\(\varphi\notin\Gamma\).
\index{consistent set!first-order}
\index{maximal consistent set!first-order}
\end{hbdefinition}

\begin{hblemma}{thm:p02-first-order-mcs-closure}{proved=annex:p2.2.1}{Closure of first-order maximal consistent sets}
If \(\Delta\subseteq\operatorname{Fm}(\Lang)\) is maximal consistent, then:
\begin{enumerate}
  \item \(\Delta\derives_{\mathsf{ND}_{\mathsf{FOL}}}\varphi\) implies
        \(\varphi\in\Delta\);
  \item exactly one of \(\varphi,\neg\varphi\) belongs to \(\Delta\);
  \item \(\varphi\land\psi\in\Delta\) iff both formulas belong to
        \(\Delta\);
  \item \(\varphi\lor\psi\in\Delta\) iff at least one belongs to
        \(\Delta\); and
  \item \(\varphi\to\psi\in\Delta\) iff
        \(\varphi\notin\Delta\) or \(\psi\in\Delta\).
\end{enumerate}
\end{hblemma}

\begin{hblemma}{thm:p02-quantifier-duality}{proved=annex:p2.2.2}{Classical quantifier duality}
The registered classical natural-deduction calculus derives the
object-language biconditionals
\[
 \neg\forall x\,\varphi\leftrightarrow\exists x\,\neg\varphi,
 \qquad
 \neg\exists x\,\varphi\leftrightarrow\forall x\,\neg\varphi.
\]
\end{hblemma}

\begin{hblemma}{thm:p02-fresh-constant}{proved=annex:p2.2.3}{Fresh-constant lemma}
Let \(c\) be a constant occurring in neither \(\Gamma\) nor \(\varphi\).
If
\(\Gamma\derives_{\mathsf{ND}_{\mathsf{FOL}}}\varphi[c/x]\), then
\(\Gamma\derives_{\mathsf{ND}_{\mathsf{FOL}}}\forall x\,\varphi\).
\end{hblemma}

\begin{hbdefinition}{def:p02-witness-property}{Witness property}
A set \(\Delta\subseteq\operatorname{Fm}(\Lang)\) has the
\emph{witness property} iff, whenever
\(\exists x\,\varphi\in\Delta\), some constant \(c\) of \(\Lang\)
satisfies \(\varphi[c/x]\in\Delta\).  This property is required because
the term model must name every existential witness.
\index{witness property}
\end{hbdefinition}

For a fresh constant \(c\), write
\(H_{x,\varphi,c}\defeq
  (\exists x\,\varphi\to\varphi[c/x])\).

\begin{hblemma}{thm:p02-witness-extension}{proved=annex:p2.2.4}{Witness-extension lemma}
If \(\Sigma\) is consistent and \(c\) occurs in neither \(\Sigma\) nor
\(\exists x\,\varphi\), then
\(\Sigma\cup\{H_{x,\varphi,c}\}\) is consistent.
\end{hblemma}

\begin{hblemma}{thm:p02-new-constant-elimination}{proved=annex:p2.2.5}{Eliminating newly added constants}
Let \(\Lang'\) extend \(\Lang\) only by constants, and let
\(\Gamma\cup\{\varphi\}\subseteq\operatorname{Fm}(\Lang)\).  If an
\(\mathsf{ND}_{\mathsf{FOL}}\)-derivation of \(\varphi\) from \(\Gamma\)
exists in \(\Lang'\), then one exists already in \(\Lang\).
\end{hblemma}

\begin{hblemma}{thm:p02-transfinite-lindenbaum}{proved=annex:p2.2.6}{Transfinite Lindenbaum lemma}
Let \(\Lang'\) be a signature, \(\lambda\) an ordinal, and
\(\langle\varphi_\alpha:\alpha<\lambda\rangle\) an enumeration of
\(\operatorname{Fm}(\Lang')\) in which every formula occurs at least once.
Every consistent \(\Sigma\subseteq\operatorname{Fm}(\Lang')\) is contained in
a maximal consistent \(\Delta\subseteq\operatorname{Fm}(\Lang')\).  Given the
enumeration, the construction uses no choice principle.
\end{hblemma}

\begin{hblemma}{thm:p02-henkin-lindenbaum}{proved=annex:p2.2.7}{Henkin--Lindenbaum lemma}
Assume the Axiom of Choice in the ambient set theory.  Every consistent
\(\Sigma\subseteq\operatorname{Fm}(\Lang)\) is contained in a maximal consistent
set \(\Delta\subseteq\operatorname{Fm}(\Lang^H)\) in a constant expansion
\(\Lang^H\), and \(\Delta\) has the witness property.
\end{hblemma}

\begin{hbmethod}{def:p02-henkin-method}{Henkin term-model method}
To turn syntactic consistency into a countermodel: expand the signature by
enough fresh constants; enumerate existential formulas and add one fresh
witness implication for each; extend the resulting consistent set by the
transfinite Lindenbaum construction; quotient terms by equality derivable
from the resulting maximal consistent witness set; interpret functions and
relations on the quotient; prove the truth lemma at the canonical assignment;
and finally take the reduct to the original signature.  The construction is
transfinite for arbitrary signatures, becomes an explicit recursion on
\(\N\) for countable signatures, and uses the Axiom of Choice only where the
entries below say it does.  Annex items P2.2.1--P2.2.12 supply the complete
method derivation.
\index{Henkin method}
\index{term-model method}
\end{hbmethod}

\section{The term structure and truth lemma}
\label{sec:p02-first-order-term-structure}

Fix a maximal consistent witness set
\(\Delta\subseteq\operatorname{Fm}(\Lang^H)\).  For terms \(r,t\), define
\[
 r\sim_\Delta t
 \quad\text{iff}\quad
 \Delta\derives_{\mathsf{ND}_{\mathsf{FOL}}}r=t.
\]
This is a relation among terms, not formula equivalence or elementary
equivalence; it is used to form the quotient below.

\begin{hblemma}{thm:p02-equality-congruence}{proved=annex:p2.2.8}{Equality congruence}
The relation \(\sim_\Delta\) is an equivalence relation.  If \(n\in\N\),
\(r_i,t_i\in\operatorname{Term}(\Lang^H)\), and
\(r_i\sim_\Delta t_i\) for \(1\leq i\leq n\), then every \(n\)-ary
function symbol \(f\) satisfies
\[
 f(r_1,\ldots,r_n)\sim_\Delta f(t_1,\ldots,t_n),
\]
and every \(n\)-ary relation symbol \(R\) satisfies
\[
 R(r_1,\ldots,r_n)\in\Delta
 \quad\text{iff}\quad
 R(t_1,\ldots,t_n)\in\Delta.
\]
\end{hblemma}

\begin{hbdefinition}{def:p02-henkin-term-structure}{Henkin term structure}
The \emph{term structure} \(\mathcal M_\Delta\) has
\[
 \dom\mathcal M_\Delta
   \defeq\operatorname{Term}(\Lang^H)/{\sim_\Delta}.
\]
Writing \([t]\) for the class of \(t\), set
\[
 f^{\mathcal M_\Delta}([t_1],\ldots,[t_n])
   \defeq[f(t_1,\ldots,t_n)]
\]
and
\[
 ([t_1],\ldots,[t_n])\in R^{\mathcal M_\Delta}
 \quad\text{iff}\quad R(t_1,\ldots,t_n)\in\Delta.
\]
Let \(s_0(x)\defeq[x]\).  Equality congruence makes the interpretations
representative-independent; the added constants make the domain nonempty.
\index{term structure!Henkin}
\end{hbdefinition}

\begin{hblemma}{thm:p02-term-model-evaluation}{proved=annex:p2.2.9}{Term evaluation in the term structure}
For every \(\Lang^H\)-term \(t\),
\[
 \operatorname{val}^{\mathcal M_\Delta}_{s_0}(t)=[t].
\]
\end{hblemma}

\begin{hblemma}{thm:p02-quantifier-membership}{proved=annex:p2.2.10}{Quantifier membership}
For every formula \(\varphi\) and variable \(x\):
\begin{enumerate}
  \item \(\exists x\,\varphi\in\Delta\) iff
        \(\varphi[t/x]\in\Delta\) for some term \(t\);
  \item \(\forall x\,\varphi\in\Delta\) iff
        \(\varphi[t/x]\in\Delta\) for every term \(t\).
\end{enumerate}
\end{hblemma}

\begin{hblemma}{thm:p02-first-order-truth-lemma}{proved=annex:p2.2.11}{First-order truth lemma}
For every \(\Lang^H\)-formula \(\varphi\),
\[
 \mathcal M_\Delta,s_0\entails\varphi
 \quad\text{iff}\quad
 \varphi\in\Delta.
\]
\end{hblemma}

\section{Completeness and compactness}
\label{sec:p02-first-order-completeness-results}

\begin{hbtheorem}{thm:p02-fol-completeness}{proved=annex:p2.2.12}{Strong completeness of classical first-order logic}
Assume the Axiom of Choice in the ambient set theory.  For every
\(\Gamma\cup\{\varphi\}\subseteq\operatorname{Fm}(\Lang)\),
\[
 \Gamma\entails\varphi
 \quad\Longrightarrow\quad
 \Gamma\derives_{\mathsf{ND}_{\mathsf{FOL}}}\varphi.
\]
Together with soundness, this identifies the registered calculus's
derivability relation with consequence over nonempty set-based structures;
it does not identify syntactic derivations with semantic truth as objects.
\end{hbtheorem}

\begin{hbcorollary}{thm:p02-fol-completeness-countable}{proved=annex:p2.2.13}{Completeness for countable signatures without choice}
If the function and relation symbols of \(\Lang\) form a countable set
(finite signatures included), then for all
\(\Gamma\cup\{\varphi\}\subseteq\operatorname{Fm}(\Lang)\),
\[
 \Gamma\entails\varphi
 \quad\Longrightarrow\quad
 \Gamma\derives_{\mathsf{ND}_{\mathsf{FOL}}}\varphi,
\]
with no use of the Axiom of Choice.
\end{hbcorollary}

\begin{hbtheorem}{thm:p02-fol-lk-completeness}{sketched=annex:p2.2.14}{Completeness of the first-order sequent calculus}
Assume the Axiom of Choice.  For finite multisets \(\Gamma,\Delta\), the
sequent \(\Gamma\Rightarrow\Delta\) is derivable in the registered
cut-admitting first-order \(\mathsf{LK}\) iff, for every \(\Lang\)-structure
\(\mathcal M\) and assignment \(s\), satisfaction of every member of
\(\operatorname{supp}(\Gamma)\) implies satisfaction of some member of
\(\operatorname{supp}(\Delta)\).  Empty antecedents and succedents are
included; no cut-free completeness claim is made.
\end{hbtheorem}

\begin{hbtheorem}{thm:p02-fol-compactness}{proved=annex:p2.2.15}{First-order compactness}
Assume the Axiom of Choice.  For every
\(T\subseteq\operatorname{Sent}(\Lang)\) and
\(\sigma\in\operatorname{Sent}(\Lang)\):
\begin{enumerate}
  \item \(T\) has a model iff every finite subset of \(T\) has a model;
  \item if \(T\entails\sigma\), then some finite \(T_0\subseteq T\)
        satisfies \(T_0\entails\sigma\).
\end{enumerate}
The two clauses are respectively the model and consequence forms.
\end{hbtheorem}

\begin{hbcorollary}{thm:p02-finite-structures-not-axiomatizable}{proved=annex:p2.2.16}{Finite structures are not first-order axiomatizable}
Assume the Axiom of Choice.  No first-order theory in the empty signature
has exactly the finite nonempty structures as its models.
\end{hbcorollary}

\section{Elementary substructures}
\label{sec:p02-first-order-elementary-substructures}

\begin{hbdefinition}{def:p02-substructure}{Substructure}
For \(\Lang\)-structures \(\mathcal N,\mathcal M\), write
\(\mathcal N\subseteq\mathcal M\) iff
\(\dom\mathcal N\subseteq\dom\mathcal M\), constants agree, the smaller
domain is closed under the functions and carries their restrictions, and
each relation of \(\mathcal N\) is the corresponding relation of
\(\mathcal M\) restricted to \(\dom\mathcal N\).
\index{substructure}
\end{hbdefinition}

\begin{hbdefinition}{def:p02-elementary-substructure}{Elementary substructure}
A substructure \(\mathcal N\subseteq\mathcal M\) is
\emph{elementary}, written \(\mathcal N\preccurlyeq\mathcal M\), iff for
every formula \(\varphi(x_1,\ldots,x_n)\) with
\(\FV(\varphi)\subseteq\{x_1,\ldots,x_n\}\) and every
\(\bar a\in(\dom\mathcal N)^n\),
\[
 \mathcal N\entails\varphi(\bar a)
 \quad\text{iff}\quad
 \mathcal M\entails\varphi(\bar a).
\]
The parameter notation abbreviates satisfaction under an assignment taking
the displayed variables to \(\bar a\).
\index{elementary substructure}
\end{hbdefinition}

\begin{hbtheorem}{thm:p02-tarski-vaught}{proved=annex:p2.2.17}{Tarski--Vaught test}
Let \(\mathcal N\subseteq\mathcal M\).  Then
\(\mathcal N\preccurlyeq\mathcal M\) iff, for every formula
\(\varphi(y,x_1,\ldots,x_n)\) with
\(\FV(\varphi)\subseteq\{y,x_1,\ldots,x_n\}\) and every
\(\bar a\in(\dom\mathcal N)^n\),
\[
 \mathcal M\entails\exists y\,\varphi(y,\bar a)
 \quad\Longrightarrow\quad
 \mathcal M\entails\varphi(b,\bar a)
 \text{ for some }b\in\dom\mathcal N.
\]
\end{hbtheorem}

\section{The L\"owenheim--Skolem theorems}
\label{sec:p02-first-order-lowenheim-skolem}

\begin{hbtheorem}{thm:p02-downward-ls}{proved=annex:p2.2.18}{Downward L\"owenheim--Skolem}
Assume the Axiom of Choice.  If \(\mathcal M\) is a \(\Lang\)-structure
and \(A\subseteq\dom\mathcal M\), then some
\(\mathcal N\preccurlyeq\mathcal M\) contains \(A\) and satisfies
\[
 \abs{\dom\mathcal N}
 \leq\max\bigl(\abs{A},\abs{\Lang},\aleph_0\bigr).
\]
\end{hbtheorem}

\begin{hbtheorem}{thm:p02-upward-ls}{proved=annex:p2.2.19}{Upward L\"owenheim--Skolem}
Assume the Axiom of Choice.  If a \(\Lang\)-theory \(T\) has an infinite
model, then for every cardinal
\(\kappa\geq\max(\abs{\Lang},\aleph_0)\), \(T\) has a model of cardinality
exactly \(\kappa\).  In particular, for every infinite \(\mathcal M\) and
such \(\kappa\), some \(\mathcal N\) of cardinality \(\kappa\) satisfies
\(\mathcal N\equiv\mathcal M\).  This does not assert that
\(\mathcal M\preccurlyeq\mathcal N\).
\end{hbtheorem}

\section{Compactness applications}
\label{sec:p02-first-order-compactness-applications}
%% Depth row D-P3-013: elementary-diagram upward LS, nonstandard models,
%% and the internal/external calibration of Skolem's paradox.

For the elementary-diagram construction and these model-theoretic
calibrations, see \cite{Hodges1993,Marker2002}.

\begin{hbdefinition}{def:p02-elementary-diagram}{Parameter expansion and elementary diagram}
For a \(\Lang\)-structure \(\mathcal M\), let
\(\Lang(\mathcal M)\) add a distinct constant \(c_a\) for each
\(a\in\dom\mathcal M\), and let \(\mathcal M_{\mathcal M}\) interpret
\(c_a\) as \(a\).  The \emph{elementary diagram}
\[
 \operatorname{ElDiag}(\mathcal M)
 \defeq
 \operatorname{Th}_{\Lang(\mathcal M)}(\mathcal M_{\mathcal M})
\]
is a theory---all parameter sentences true in that expansion---not a
structure.  Unlike the parameter-free theory
\(\operatorname{Th}_{\Lang}(\mathcal M)\), it records the truth of every
formula at each named tuple; that extra information is exactly what the
elementary-extension construction uses.
\index{elementary diagram}
\index{parameter expansion}
\end{hbdefinition}

\begin{hbtheorem}{thm:p02-upward-ls-elementary-extension}{proved=annex:p2.2.20}{Upward L\"owenheim--Skolem via elementary diagrams}
Assume the Axiom of Choice.  Let \(\mathcal M\) be an infinite
\(\Lang\)-structure and
\[
 \kappa\geq
 \max\bigl(\abs{\Lang},\abs{\dom\mathcal M},\aleph_0\bigr).
\]
Then there are a \(\Lang\)-structure \(\mathcal N\) of cardinality exactly
\(\kappa\) and an elementary embedding
\(j:\mathcal M\to\mathcal N\) whose range is proper.
\end{hbtheorem}

\begin{hbcorollary}{thm:p02-nonstandard-arithmetic-model}{proved=annex:p2.2.21}{A countable nonstandard elementary model of arithmetic}
Assume the Axiom of Choice.  Then the standard first-order structure
\(\mathcal N_{\mathrm{ar}}=(\N,0,S,+,\cdot)\) has a countable proper
elementary extension.  Consequently a countable model satisfies every
first-order sentence true in \(\mathcal N_{\mathrm{ar}}\) while containing
elements outside the embedded standard natural numbers.
\end{hbcorollary}

\begin{hbwarning}{def:p02-skolem-paradox}{Skolem's paradox is an internal--external distinction}
Assume the Axiom of Choice.  Conditionally, if a first-order set theory in a
countable signature has an infinite model, downward L\"owenheim--Skolem
supplies an externally countable elementary model.  Such a model can satisfy
that one of its sets is
\emph{uncountable}: internally there is no bijection belonging to the model,
while an external enumeration of the model need not be one of its elements.
There is no contradiction and no unconditional model-existence claim here;
the quantifiers range over different domains.  See \cite{Skolem1934,Hodges1993}.
\index{Skolem's paradox}
\end{hbwarning}
